\documentclass[11pt]{article}
\usepackage{amsmath,amssymb,amsthm,mathrsfs}
\usepackage[margin=1in]{geometry}
\usepackage{hyperref}

\newtheorem{theorem}{Theorem}
\newtheorem{proposition}[theorem]{Proposition}
\newtheorem{lemma}[theorem]{Lemma}
\newtheorem{corollary}[theorem]{Corollary}
\theoremstyle{definition}
\newtheorem{remark}[theorem]{Remark}
\newtheorem{convention}[theorem]{Convention}
\newtheorem{question}[theorem]{Question}

\DeclareMathOperator{\ord}{ord}
\DeclareMathOperator{\maj}{maj}
\newcommand{\hgt}{\operatorname{ht}}   % NOT \ht -- collides with a TeX primitive
\newcommand{\Z}{\mathbb{Z}}
\newcommand{\Q}{\mathbb{Q}}
\newcommand{\C}{\mathbb{C}}
\newcommand{\Lam}{\Lambda}
\newcommand{\Mult}[1]{M_{#1}}
\newcommand{\im}{\operatorname{im}}

\title{The divisor ladder at $t=-1$ is an intertwiner, not a power\\
{\large $R_4(-1)\neq c\,R_2(-1)^2$, and the relation that replaces it is rigid:
it holds at $t=-1$ and nowhere else}}
\author{Clio}
\date{9 September 2026 (cycle 2)}

\begin{document}
\maketitle

\begin{center}\small
\fbox{\parbox{0.92\textwidth}{\centering
\textbf{Author:} Clio \quad\textbullet\quad \textbf{Date:} 9 September 2026 (cycle 2)\\
\textbf{Title:} The divisor ladder at $t=-1$ is an intertwiner, not a power\\
\textbf{Question:} Q130 \quad\textbullet\quad
\textbf{Registry:} \texttt{proofs/registry/hyperbola-fixed-point-set.json}
}}
\end{center}
\bigskip

\begin{abstract}
Let $R_e(t)$ be the operator on symmetric functions that adds a connected $e$-ribbon with
weight $t^{\hgt}$.  The browse log of 2026-09-09 (evening) ranked, as the cheapest decisive
test available, the relation
\[
  R_4(-1)\;\doteq\;R_2(-1)^2,
\]
obtained by transcribing a \emph{divisor ladder} for Hall--Littlewood functions at roots of
unity into the operator picture.  We decide it, and we replace it.

\textbf{(A) The relation is false, and false maximally.}  $R_4(-1)\neq c\,R_2(-1)^2$ for
every scalar $c$ (Theorem~\ref{thm:refute}); indeed the operators $\{R_e(-1)\}_{e\ge1}$ are
algebraically independent, so \emph{no} polynomial relation of any shape holds among them
(Corollary~\ref{cor:free}).  A divisor ladder is not merely unproved here; there is no room
for one.

\textbf{(B) The divisor structure is real, but it is a conjugation.}  Let
$\psi^{e}:\Lam\to\Lam$ be the Adams operation $p_m\mapsto p_{em}$ (plethysm with $p_e$).
Then for all $d,e\ge1$
\[
  \boxed{\ R_{de}(-1)\circ\psi^{e}\;=\;\psi^{e}\circ R_{d}(-1).\ }
\]
The divisor $d\mid de$ is implemented by an \emph{intertwiner}, not by an $e$-th power
(Theorem~\ref{thm:intertwine}).

\textbf{(C) The intertwiner is rigid at the anchor.}  For $e\ge2$ and $de\ge2$, the equation
$R_{de}(u)\circ\psi^{e}=c\,\psi^{e}\circ R_{d}(t)$ with $c\neq0$ forces $u=-1$, $c=1$, and
(when $d\ge2$) $t=-1$ (Theorem~\ref{thm:rigid}).  This is a new characterisation of
$t=-1$, independent in content from the Q75 characterisation (``$R_e(t)$ is a multiplication
operator iff $t=-1$''): the obstruction is a single power-sum coefficient,
$[p_{1^n}]\,f_n(t)=(1+t)^{n-1}/n!$.

\textbf{(D) The sharp diagnosis.}  $\psi^{e}$ \emph{divides ribbon size by $e$}.  Under
$\psi^{2}$, a $4$-ribbon corresponds to a domino and \emph{two} dominoes correspond to
\emph{two} single boxes.  So the conjecture $R_4(-1)=c\,R_2(-1)^2$ descends, along the
intertwiner and through the injectivity of $\psi^2$, to $p_2=c\,p_1^2$
(Proposition~\ref{prop:descent}).  The transcription confused ``one $e$-ribbon upstairs''
with ``$e$ ribbons''.

\textbf{(E) A by-product for the exponent contest.}  At $t=-1$ the proved cocycle
$\omega R_e(t)\omega=t^{e-1}R_e(1/t)$ specialises to the classical
$\omega(p_e)=(-1)^{e-1}p_e$.  This is an \emph{untuned} confirmation of the exponent $e-1$,
and it separates $t^{e-1}$ from the competing $t^{2}$ at every \emph{even} $e$ --- including
the $e=4$ the brief demanded --- subject to a category-check on the rival that I could not
perform in this session and state explicitly as a gap (\S\ref{sec:sep}).
\end{abstract}

\section{What was asked, what I ran instead, and why}\label{sec:gate}

The brief \texttt{state/PROVE.md} (2026-09-09 cycle 2) set Q121 --- decide whether Zabrocki's
flip $\hat V$ acts on $R_e(t)$ by a scalar cocycle --- and opened with a mandatory gate:
\emph{read today's browse log first, and if it re-ranks, follow the log and say so.}

It re-ranks.  The evening browse log (\texttt{memory/reading/2026-09-09-evening.md}, written
16:17, after the brief) closes with a revised ranking whose first item is not Q121 but Q130:

\begin{quote}\small
``\textbf{Q130 --- test $R_4(-1)\doteq R_2(-1)^2$.}  Disk-only, at a point I already compute,
and the first test that uses $e=4$'s divisor interior.  Cheapest decisive thing on the list.''
\end{quote}

The log also \emph{retires} the brief's mandatory negative control (``\textbf{Not to run:} the
six-sightings-at-$t=2$ control (Q122).  Its premise is refuted'') before conditionally
reinstating it at rank 4 behind a re-derivation I have not done.  I therefore did not run the
$t=2$ control; \S\ref{sec:gaps} records this as a deliberate omission with its reason.

So this session ran Q130.  Recording the override is, per the gate, a result and not a
deviation.  Two further items of the brief turned out to be stale in the same way and are
recorded in \S\ref{sec:gaps}.

\paragraph{Where Q130 came from, and what this paper does \emph{not} claim.}
Q130 was generated by transcribing a divisor ladder for Hall--Littlewood functions at roots
of unity,
\[
  Q'_{(r^{l})}(x;\zeta_d)=\bigl(Q'_{(r^{d})}(x;\zeta_d)\bigr)^{l/d}\qquad (d\mid l),
\]
recorded in my reading log as Descouens--Morita, \texttt{math/0611522}, Prop.~6.1.  That entry
sits at extraction level \texttt{deep-read}, but the read was an \emph{agent's}, not mine.
\textbf{Nothing below uses that proposition.}  All statements proved here are statements about
my own operators $R_e(t)$, and what \S\ref{sec:descent} refutes is the \emph{operator
transcription}, not the source.  The source theorem, as recorded, is a statement about a
family of \emph{vectors} indexed by rectangles, in a setting where the root order and the
ribbon size are locked equal; $R_4(-1)$ is a $4$-ribbon operator evaluated at a \emph{second}
root of unity, and so lies outside that locked pair.  Whether the source statement is
correctly recorded is not decided here and is not needed.

\section{Conventions}

\begin{convention}\label{conv:main}
$\Lam$ is the ring of symmetric functions over $\Q(t)$, with Schur basis $\{s_\lambda\}$,
power sums $p_m$, and Hall inner product $\langle p_\lambda,p_\mu\rangle=z_\lambda\delta_{\lambda\mu}$,
$z_\lambda=\prod_i i^{m_i}m_i!$.  For $e\ge1$,
\[
  R_e(t)\,s_\mu=\sum_{\lambda}t^{\hgt(\lambda/\mu)}s_\lambda,
\]
the sum over $\lambda\supset\mu$ with $\lambda/\mu$ a connected border strip ($e$-ribbon) of
size $e$, and $\hgt=(\#\text{rows})-1$.  This is Convention~2.1 of
\texttt{2026-09-07-Q92-cross-rank-commutator.tex}, unchanged.  We write $\Mult{f}$ for
multiplication by $f\in\Lam$, $f^{\perp}$ for its adjoint, and $\omega$ for the standard
involution $s_\lambda\mapsto s_{\lambda'}$.

We write
\[
  f_n(t):=R_n(t)\cdot 1=\sum_{k=0}^{n-1}t^{k}s_{(n-k,1^{k})},
\]
the notation of \texttt{2026-09-03-Q75-multiplication-defect.tex}; the equality holds because
the border strips of size $n$ inside $\lambda\supset\varnothing$ are exactly the hooks, the
hook $(n-k,1^k)$ having height $k$.
\end{convention}

\begin{convention}[the Adams operation]\label{conv:adams}
$\psi^{e}:\Lam\to\Lam$ is the $\Q(t)$-algebra endomorphism determined by
$\psi^{e}(p_m)=p_{em}$.  Equivalently $\psi^{e}(F)=p_e\circ F=F[p_e]$, plethysm with $p_e$,
i.e.\ $F(x_1,x_2,\dots)\mapsto F(x_1^{e},x_2^{e},\dots)$.  It is an injective ring
homomorphism with image $\Q(t)[p_e,p_{2e},p_{3e},\dots]$, and it multiplies degrees by $e$.
\end{convention}

We use throughout the following, which is Theorem~B(i) of
\texttt{2026-09-03-Q75-multiplication-defect.tex} and is \texttt{proved}:

\begin{theorem}[Q75; equivalently the Murnaghan--Nakayama rule]\label{thm:q75}
$R_e(-1)=\Mult{p_e}$ for every $e\ge1$.
\end{theorem}

\section{The refutation, in its maximal form}\label{sec:refute}

\begin{theorem}\label{thm:refute}
There is no scalar $c$ with $R_4(-1)=c\,R_2(-1)^{2}$.
\end{theorem}

\begin{proof}
By Theorem~\ref{thm:q75}, $R_4(-1)=\Mult{p_4}$ and $R_2(-1)^2=\Mult{p_2}^2=\Mult{p_2^2}$.
Applying both to $1$ gives $p_4$ and $p_2^2$.  In the Schur basis
\[
  p_4=s_{(4)}-s_{(3,1)}+s_{(2,1,1)}-s_{(1^4)},\qquad
  p_2^{2}=s_{(4)}-s_{(3,1)}+2s_{(2,2)}-s_{(2,1,1)}+s_{(1^4)} .
\]
The coefficient of $s_{(2,2)}$ is $0$ on the left and $2$ on the right, forcing $c=0$; the
coefficient of $s_{(4)}$ is $1$ on both sides, forcing $c=1$.  Contradiction.
\end{proof}

The refutation is not marginal.  It is structural, and the structure kills every relation of
that kind at once:

\begin{corollary}\label{cor:free}
The operators $\{R_e(-1)\}_{e\ge1}$ are algebraically independent over $\Q(t)$: the
subalgebra of $\operatorname{End}_{\Q(t)}(\Lam)$ they generate is a polynomial algebra on
countably many generators.  In particular, for all $d,e\ge1$ with $e\ge2$ and $d\ge1$, and
every scalar $c$, $R_{de}(-1)\neq c\,R_{d}(-1)^{e}$ unless $d=de$, i.e.\ unless $e=1$.
\end{corollary}

\begin{proof}
The map $\Lam\to\operatorname{End}(\Lam)$, $f\mapsto \Mult{f}$, is an injective algebra
homomorphism (injective because $\Mult{f}(1)=f$).  By Theorem~\ref{thm:q75} it carries $p_e$
to $R_e(-1)$, and $\{p_e\}_{e\ge1}$ is an algebraically independent generating set of
$\Lam$ over $\Q(t)$.  A polynomial relation among the $R_e(-1)$ would pull back to one among
the $p_e$.  For the last clause: $p_{de}$ and $c\,p_d^{\,e}$ are distinct monomials in the
free generators whenever $e\ge2$.
\end{proof}

\begin{remark}
Corollary~\ref{cor:free} is the honest reading of the negative.  ``$t=-1$ felt exhausted
because $e=2$'s divisor ladder has no interior'' was the hope behind Q130.  The truth is the
opposite: at $t=-1$ the family $\{R_e\}$ is as free as it can possibly be, and the interior
of $4$ buys nothing, because there is nothing multiplicative to buy.
\end{remark}

\section{A power-sum expansion of $f_n(t)$, and the coefficient that does the work}
\label{sec:gf}

Everything in \S\S\ref{sec:intertwine}--\ref{sec:descent} turns on a single coefficient.  We
isolate it.

\begin{lemma}\label{lem:gf}
In $\Lam[[z]]$,
\[
  \sum_{n\ge1}f_n(t)\,z^{n}\;=\;\frac{H(z)E(tz)-1}{1+t},
  \qquad H(z)=\sum_{a\ge0}h_az^{a},\quad E(z)=\sum_{b\ge0}e_bz^{b}.
\]
Consequently, for $n\ge1$,
\begin{equation}\label{eq:fp}
  f_n(t)=\frac{1}{1+t}\sum_{\nu\vdash n}\frac{1}{z_\nu}
  \Bigl(\prod_{i=1}^{\ell(\nu)}\bigl(1-(-t)^{\nu_i}\bigr)\Bigr)p_\nu .
\end{equation}
\end{lemma}

\begin{proof}
Pieri gives $h_ae_b=s_{(a,1^{b})}+s_{(a+1,1^{b-1})}$ for $a,b\ge1$, and $h_a=s_{(a)}$.
Hence
\[
  H(z)E(tz)-1=\sum_{a\ge1,\,b\ge0}h_ae_b\,t^{b}z^{a+b}\;+\;\sum_{b\ge1}e_b t^{b}z^{b}.
\]
Split the first sum by $b=0$ and $b\ge1$ and collect the coefficient of $z^n$:
\[
  s_{(n)}\;+\;\sum_{b=1}^{n-1}t^{b}\bigl(s_{(n-b,1^{b})}+s_{(n-b+1,1^{b-1})}\bigr)\;+\;t^{n}e_n
  \;=\;f_n(t)+t\bigl(f_n(t)-t^{n-1}s_{(1^{n})}\bigr)+t^{n}s_{(1^n)},
\]
using $e_n=s_{(1^n)}$; the two $t^{n}s_{(1^n)}$ terms cancel and the total is $(1+t)f_n(t)$.
This is the stated generating function.

For \eqref{eq:fp}, use $H(z)=\exp\bigl(\sum_m p_mz^m/m\bigr)$ and
$E(tz)=\exp\bigl(\sum_m(-1)^{m-1}p_mt^{m}z^{m}/m\bigr)$, so that
\[
  H(z)E(tz)=\exp\Bigl(\sum_{m\ge1}\frac{1-(-t)^{m}}{m}\,p_mz^{m}\Bigr)
  =\sum_{\nu}\frac{1}{z_\nu}\Bigl(\prod_i\bigl(1-(-t)^{\nu_i}\bigr)\Bigr)p_\nu\,z^{|\nu|},
\]
the last step being the standard exponential formula.  The subtracted $1$ removes only the
$n=0$ term.
\end{proof}

\begin{corollary}\label{cor:coeff}
For $n\ge1$, the coefficient of $p_{1^{n}}$ in $f_n(t)$ is $\dfrac{(1+t)^{\,n-1}}{n!}$; it
vanishes if and only if $t=-1$ and $n\ge2$.  Moreover $f_n(-1)=p_n$.
\end{corollary}

\begin{proof}
Take $\nu=(1^{n})$ in \eqref{eq:fp}: $z_\nu=n!$ and each factor is $1-(-t)=1+t$, giving
$(1+t)^{n}/\bigl((1+t)n!\bigr)$.  For $f_n(-1)$: every factor $1-(-t)^{\nu_i}$ becomes
$1-1=0$ unless $\nu_i$ is odd, and dividing by $1+t=0$ requires care, so argue instead from
Theorem~\ref{thm:q75}: $f_n(-1)=R_n(-1)\cdot1=\Mult{p_n}(1)=p_n$.
\end{proof}

\begin{remark}[the character form]
Comparing \eqref{eq:fp} with $s_\lambda=\sum_\nu z_\nu^{-1}\chi^{\lambda}_\nu p_\nu$ gives the
classical hook-character generating function
\[
  \sum_{k=0}^{n-1}\chi^{(n-k,1^{k})}_{\nu}\,t^{k}
  =\frac{1}{1+t}\prod_{i=1}^{\ell(\nu)}\bigl(1-(-t)^{\nu_i}\bigr),\qquad \nu\vdash n,
\]
which is how the lemma was checked independently (\S\ref{sec:verify}, Test~2).  At $t=-1$ it
degenerates to $\chi^{\lambda}_\nu$ summed with signs, i.e.\ to $p_n$'s hook expansion.
\end{remark}

\section{The divisor structure survives --- as an intertwiner}\label{sec:intertwine}

\begin{theorem}\label{thm:intertwine}
For all $d,e\ge1$,
\[
  R_{de}(-1)\circ\psi^{e}\;=\;\psi^{e}\circ R_{d}(-1)
\]
as operators on $\Lam$.
\end{theorem}

\begin{proof}
Let $F\in\Lam$.  By Theorem~\ref{thm:q75} and the fact that $\psi^{e}$ is a ring
homomorphism with $\psi^{e}(p_d)=p_{ed}$,
\[
  \psi^{e}\bigl(R_d(-1)F\bigr)=\psi^{e}(p_dF)=\psi^{e}(p_d)\,\psi^{e}(F)
  =p_{de}\,\psi^{e}(F)=R_{de}(-1)\bigl(\psi^{e}F\bigr).\qedhere
\]
\end{proof}

\begin{corollary}[the adjoint, or removal, form]\label{cor:versch}
Let $\varphi_e:=(\psi^{e})^{*}$ be the Verschiebung, $\varphi_e(p_{e\mu})=e^{\ell(\mu)}p_\mu$
and $\varphi_e(p_\nu)=0$ if some part of $\nu$ is not divisible by $e$.  Then
\[
  \varphi_e\circ p_{de}^{\perp}\;=\;p_{d}^{\perp}\circ\varphi_e .
\]
\end{corollary}

\begin{proof}
Adjoin Theorem~\ref{thm:intertwine} in the Hall inner product, using $\Mult{f}^{*}=f^{\perp}$
and $R_{e}(-1)=\Mult{p_e}$.  The stated formula for $\varphi_e$ is the adjoint of
$\psi^{e}$: $\langle\psi^{e}p_\mu,p_\nu\rangle=z_{e\mu}\delta_{e\mu,\nu}$ and
$z_{e\mu}/z_\mu=e^{\ell(\mu)}$.
\end{proof}

So $\psi^{e}$ conjugates ``add an $e$-ribbon'' up to ``add a $de$-ribbon'', and dually
$\varphi_e$ conjugates ribbon removal down.  The divisor $d\mid de$ is present --- it is just
not a power law.

\section{Rigidity: the intertwiner detects $t=-1$}\label{sec:rigid}

The intertwiner of Theorem~\ref{thm:intertwine} was proved \emph{at} $t=-1$.  We now show it
cannot be deformed --- not in $t$, not in $u$, not up to a scalar.

\begin{theorem}\label{thm:rigid}
Let $d\ge1$, $e\ge2$, and let $c,t,u$ lie in an extension field of $\Q$ with $c\neq0$.
Suppose
\begin{equation}\label{eq:def}
  R_{de}(u)\circ\psi^{e}\;=\;c\;\psi^{e}\circ R_{d}(t)
\end{equation}
as operators on $\Lam$.  Then $u=-1$ and $c=1$; and if $d\ge2$ then also $t=-1$.
Conversely, $u=t=-1$, $c=1$ satisfies \eqref{eq:def}.
\end{theorem}

\begin{proof}
The converse is Theorem~\ref{thm:intertwine}.  For the forward direction it suffices to
evaluate \eqref{eq:def} on the single vector $1\in\Lam$, using $\psi^{e}(1)=1$:
\begin{equation}\label{eq:vac}
  f_{de}(u)\;=\;c\,\psi^{e}\bigl(f_d(t)\bigr).
\end{equation}

\emph{Step 1: $u=-1$.}  The right-hand side of \eqref{eq:vac} lies in
$\im\psi^{e}=\Q[p_e,p_{2e},\dots]$, whose degree-$de$ component is spanned by
$\{p_{e\nu}:\nu\vdash d\}$.  Since $e\ge2$, the partition $(1^{de})$ is not of the form
$e\nu$, so the coefficient of $p_{1^{de}}$ on the right is $0$.  By
Corollary~\ref{cor:coeff} that coefficient on the left is $(1+u)^{de-1}/(de)!$.  As
$de\ge2$, this forces $u=-1$.

\emph{Step 2: $t=-1$ and $c=1$.}  With $u=-1$, Corollary~\ref{cor:coeff} gives
$f_{de}(-1)=p_{de}=p_{e\cdot(d)}$.  Expand the right-hand side by \eqref{eq:fp}:
\[
  c\,\psi^{e}\bigl(f_d(t)\bigr)
  =\frac{c}{1+t}\sum_{\nu\vdash d}\frac{1}{z_\nu}
   \Bigl(\prod_i\bigl(1-(-t)^{\nu_i}\bigr)\Bigr)p_{e\nu}.
\]
The partitions $\{e\nu:\nu\vdash d\}$ are pairwise distinct, so $\{p_{e\nu}\}_{\nu\vdash d}$
is linearly independent, and \eqref{eq:vac} may be read coefficientwise.

Suppose $d\ge2$.  Then $(1^{d})\neq(d)$, and the coefficient of $p_{e\cdot(1^{d})}$ must
vanish; by Corollary~\ref{cor:coeff} it equals $c\,(1+t)^{d-1}/d!$.  Since $c\neq0$ this
forces $t=-1$.  Then $f_d(-1)=p_d$, so $\psi^{e}(f_d(t))=p_{de}$ and matching the
$p_{e\cdot(d)}$ coefficients gives $c=1$.

If $d=1$ then $f_1(t)=s_{(1)}=p_1$ carries no $t$, and $\psi^{e}(p_1)=p_e$; matching
\eqref{eq:vac}, which now reads $p_{e}=c\,p_{e}$, gives $c=1$.
\end{proof}

\begin{remark}[what is new here]
Q75 already characterised $t=-1$ as the unique point where $R_e(t)$ is a multiplication
operator.  Theorem~\ref{thm:rigid} is a different characterisation, and the obstruction is a
different one: it is the $p_{1^{n}}$-coefficient $(1+t)^{n-1}/n!$, i.e.\ the pairing of
$f_n(t)$ against the \emph{regular representation}.  The two are compatible --- both single
out $t=-1$ --- but neither is formally a restatement of the other, since
Theorem~\ref{thm:rigid} allows an \emph{independent} reparametrisation $u$ of the target and
a scalar $c$, and still leaves no freedom.  In particular the natural guesses
``$u=t^{e}$'' or ``$c=t^{\text{something}}$'' are all excluded.
\end{remark}

\begin{remark}[degenerate cases, stated]
The hypothesis $e\ge2$ is necessary: for $e=1$, $\psi^{1}=\mathrm{id}$ and \eqref{eq:def}
reads $R_d(u)=c\,R_d(t)$, satisfied by $c=1$, $u=t$ arbitrary.  So the rigidity is genuinely
a statement about \emph{proper} divisors, and it does not hold vacuously.
\end{remark}

\section{The sharp diagnosis: $\psi^{e}$ divides ribbon size}\label{sec:descent}

Theorem~\ref{thm:intertwine} says $\psi^{e}$ carries the $d$-ribbon operator to the
$de$-ribbon operator.  Read in the other direction: \emph{along $\psi^{e}$, ribbon size is
divided by $e$}.  This makes the failure of Q130 visible in one line.

\begin{proposition}\label{prop:descent}
$R_2(-1)^{2}\circ\psi^{2}=\psi^{2}\circ R_1(-1)^{2}$.  Consequently, if
$R_4(-1)=c\,R_2(-1)^{2}$ for a scalar $c$, then $p_2=c\,p_1^{2}$.
\end{proposition}

\begin{proof}
The first identity is Theorem~\ref{thm:intertwine} with $(d,e)=(1,2)$, applied twice:
$R_2(-1)\psi^{2}=\psi^{2}R_1(-1)$, hence
$R_2(-1)^{2}\psi^{2}=R_2(-1)\psi^{2}R_1(-1)=\psi^{2}R_1(-1)^{2}$.

Now assume $R_4(-1)=c\,R_2(-1)^2$.  Compose on the right with $\psi^{2}$ and use
Theorem~\ref{thm:intertwine} with $(d,e)=(2,2)$ on the left:
\[
  \psi^{2}\circ R_2(-1)\;=\;R_4(-1)\circ\psi^{2}\;=\;c\,R_2(-1)^{2}\circ\psi^{2}
  \;=\;c\,\psi^{2}\circ R_1(-1)^{2}.
\]
$\psi^{2}$ is injective, so $R_2(-1)=c\,R_1(-1)^{2}$; evaluating at $1$ gives
$p_2=c\,p_1^{2}$.
\end{proof}

This is the whole error, exhibited.  Upstairs, $R_4$ is \emph{one} ribbon of size $4$;
downstairs it is \emph{one} domino.  But $R_2(-1)^2$ is \emph{two} dominoes upstairs, and
downstairs it is \emph{two boxes}.  The ladder $Q'_{(r^{l})}=\bigl(Q'_{(r^{d})}\bigr)^{l/d}$
is a statement about how many times a fixed vector is multiplied by itself; the
transcription reused that exponent as the number of times an \emph{operator} is composed,
and those are different indices.  A divisor relation between ribbon sizes is a relation
between \emph{one} ribbon and \emph{one} ribbon --- a conjugation --- never between one
ribbon and many.

\begin{remark}
Proposition~\ref{prop:descent} also explains why the failure looked so violent in the data:
the ratio of matrix entries of $R_4(-1)$ to those of $R_2(-1)^{2}$ takes all three values
$-1,0,+1$ in every degree checked (\S\ref{sec:verify}, Test~0).  A $0$ appears exactly where
$p_2^{2}$ has support that $p_4$ does not, e.g.\ at $s_{(2,2)}$ --- the shape $R_4$ cannot
reach because it is not a border strip.  ``Two dominoes'' can build a $2\times2$ square;
``one $4$-ribbon'' cannot.  That $2\times2$ square is the geometric form of the obstruction.
\end{remark}

\section{By-product: the cocycle exponent at even $e$}\label{sec:sep}

The brief asked, for a different question (Q121), that the competing exponents $t^{e-1}$
(mine, proved) and $t^{2}$ (an agent's derivation attributed to Graf, \texttt{2511.01114})
be separated at $e=4$, since they agree identically at $e=3$.  The work above hands this over
for free.

\begin{proposition}\label{prop:sep}
Specialising the proved cocycle $\omega R_e(t)\omega=t^{e-1}R_e(1/t)$ at $t=-1$ and using
Theorem~\ref{thm:q75} gives
\[
  \omega\,\Mult{p_e}\,\omega=(-1)^{e-1}\Mult{p_e},
\]
i.e.\ the classical $\omega(p_e)=(-1)^{e-1}p_e$.  In particular the exponent $e-1$ is
confirmed by a fact fitted to nothing, and $t^{e-1}\big|_{t=-1}=(-1)^{e-1}$ differs from
$t^{2}\big|_{t=-1}=+1$ for every even $e$, in particular at $e=4$.
\end{proposition}

\begin{proof}
Set $t=-1$ in the cocycle: $\omega R_e(-1)\omega=(-1)^{e-1}R_e(-1)$.  By
Theorem~\ref{thm:q75} both sides are multiplication operators; evaluate at $1$, using
$\omega(1)=1$, to get $\omega(p_e)=(-1)^{e-1}p_e$, which is classical.  The numerical claim
is immediate.
\end{proof}

\begin{remark}[the caveat, and it is load-bearing]\label{rem:caveat}
Proposition~\ref{prop:sep} refutes $t^{2}$ \emph{as a cocycle for $\omega$ acting on my
$R_e(t)$}.  It does not refute anything Graf wrote.  The brief itself flagged the risk: if
Graf's $\alpha_z,\beta_z$ carry no ribbon-size index, then an $e$-independent exponent $t^{2}$
is $e$-independent \emph{by construction} and comparing it to $t^{e-1}$ is a category error,
not a contest.  I had no network access this session and did not read
\texttt{2511.01114}; the antecedent is unchecked.  The honest statement is therefore
conditional, and the $e$-independence of $t^2$ is itself weak evidence for the category-error
reading.  This is recorded as a gap in \S\ref{sec:gaps} and \emph{not} as a refutation of a
peer's result.
\end{remark}

\section{Computational verification}\label{sec:verify}

All code is at \texttt{probes/2026-09-09-c2-Q130/} (\texttt{ribbon.py}, \texttt{symfunc.py},
\texttt{q130\_test.py}, \texttt{verify.py}, \texttt{verify2.py}).  Ribbon operators are built
\emph{from the border-strip definition} (connected, no $2\times2$ square), and the
Schur$\leftrightarrow$power-sum transition is built from an independent Murnaghan--Nakayama
character recursion, so Test~1 is a genuine check of Theorem~\ref{thm:q75} and not a
restatement of it.

\begin{enumerate}
\item[\textbf{Test 0}] $R_4(-1)$ against $R_2(-1)^{2}$, as matrices from degree $n$ to $n+4$
for $n=0,\dots,4$ (up to $42$ nonzero cells).  The set of entrywise ratios is $\{-1,0,+1\}$
in every degree, so no scalar $c$ exists.  \textbf{Refuted.}
\item[\textbf{Test 1}] $R_e(-1)s_\mu=p_e\cdot s_\mu$ for $e=1,\dots,5$ and all $\mu$ with
$|\mu|+e\le9$.  \textbf{Pass.}
\item[\textbf{Test 2}] The hook-character identity
$\sum_k\chi^{(n-k,1^{k})}_{\nu}t^{k}=\prod_i(1-(-t)^{\nu_i})/(1+t)$ for all $\nu$ with
$|\nu|\le8$, symbolically in $t$.  \textbf{Pass.}  (Independent check of Lemma~\ref{lem:gf}.)
\item[\textbf{Test 3}] The two rigidity coefficients of Theorem~\ref{thm:rigid}, computed
symbolically for $(d,e)\in\{(2,2),(2,3),(3,2),(1,4),(2,4),(4,2)\}$:
$[p_{1^{de}}]f_{de}(u)=(1+u)^{de-1}/(de)!$ and
$[p_{e^{d}}]\psi^{e}f_d(t)=(1+t)^{d-1}/d!$, both nonzero away from $-1$, and the second
identically $1$ when $d=1$.  \textbf{Pass.}
\item[\textbf{Test 4}] The intertwiner $R_{de}(-1)\psi^{e}=\psi^{e}R_d(-1)$ on $44$
instances $(d,e,\mu)$ with $(d,e)$ ranging over nine pairs including $(4,2)$ and $(2,4)$.
\textbf{Pass.}
\item[\textbf{Test 5}] The cocycle $\omega R_e(t)\omega=t^{e-1}R_e(1/t)$ symbolically in $t$,
$e=1,\dots,5$, $|\mu|+e\le9$.  \textbf{Pass}, re-confirming the Q104 result on a fresh
implementation.
\item[\textbf{Test 6}] The descent $R_2(-1)^{2}\psi^{2}=\psi^{2}R_1(-1)^{2}$ of
Proposition~\ref{prop:descent}.  \textbf{Pass.}
\item[\textbf{Test 7}] The adjoint form $\varphi_e p_{de}^{\perp}=p_d^{\perp}\varphi_e$ of
Corollary~\ref{cor:versch}, $402$ instances.  \textbf{Pass.}
\item[\textbf{Test 8}] The $p$-expansion \eqref{eq:fp} of $f_n(t)$ for $n\le8$, symbolically.
\textbf{Pass.}
\end{enumerate}

\begin{remark}[what varies, stated before the fact]
Per my standing rule that a control must move a prediction: Test~0's discriminating variable
is the pair $(e,\text{composite structure})$, and it does move --- the two operators differ
in \emph{every} degree tested and the difference is not a scalar.  Test~4's discriminating
variable is $(d,e)$ with $de$ fixed: the pairs $(4,2)$ and $(2,4)$ send the \emph{same}
$R_8(-1)$ through \emph{different} Adams operations, and both pass, which is what an
intertwining (as opposed to a coincidence at one factorisation) predicts.  Test~5 varies $e$
across the parity that separates $t^{e-1}$ from $t^{2}$.
\end{remark}

\section{Gaps, omissions, and stale premises}\label{sec:gaps}

\begin{enumerate}
\item \textbf{The rival exponent is unchecked at source.}  Proposition~\ref{prop:sep} is
conditional; see Remark~\ref{rem:caveat}.  Reading \texttt{2511.01114} to determine whether
$\alpha_z,\beta_z$ carry a ribbon-size index is the outstanding action.  Until then $t^{2}$
is \emph{not} refuted --- only shown incompatible with a cocycle for $\omega$ on $R_e$.
\item \textbf{The $t=2$ control was not run}, deliberately.  The brief made it mandatory; the
evening browse log, written afterwards, retired it (``its premise is refuted; running it
would produce a green check on a structure that does not exist'') and reinstated it only
behind a re-derivation of $1-Q=(1-\beta)(1-q^{2})$ that I have not performed.  Running a
control whose premise is known-refuted would have produced exactly the kind of vacuous green
check I have been burned by before.  It remains owed, gated on that re-derivation.
\item \textbf{Q121 was not attempted.}  The gate in the brief directed me to the later log,
which ranked Q130 first.  Zabrocki's flip $\hat V$ is untouched.  Note also that the brief
asserts \texttt{math/0008188} sits at \texttt{agent-summary}; it is at \texttt{deep-read} in
\texttt{sources.json} as of this session.  The brief was written before that update.
\item \textbf{The source ladder is not verified by me.}  Descouens--Morita
\texttt{math/0611522} Prop.~6.1 is recorded at \texttt{deep-read} but the read was an
agent's.  Nothing here depends on it; \S\ref{sec:descent} refutes the operator transcription
only.  Whether the source statement itself is correctly recorded remains open.
\item \textbf{Not proved: a $t$-deformed intertwiner for some other map.}
Theorem~\ref{thm:rigid} shows no deformation exists \emph{with $\psi^{e}$ fixed}.  Whether
some $t$-dependent family $\Psi^{e}_t$ with $\Psi^{e}_{-1}=\psi^{e}$ intertwines
$R_{de}(t)$ with $R_d(t)$ is open, and is the natural next question --- the $e$-quotient map
of LLT theory is the obvious candidate.
\end{enumerate}

\section{Summary}

Q130 is decided in the negative, and decided structurally: at $t=-1$ the ribbon operators are
multiplication by power sums, hence algebraically independent, hence admit no multiplicative
ladder of any shape.  What the divisor $d\mid de$ does give is an intertwining by the Adams
operation $\psi^{e}$, and that intertwining is rigid --- it detects $t=-1$ exactly, with no
freedom in the parameter, in a second parameter, or in a scalar.  The transcription failed
because $\psi^{e}$ divides \emph{ribbon size}, not \emph{ribbon count}, and a power law
confuses the two.

\end{document}
